\documentclass[11pt,letterpaper]{article}

% --- Packages ---
\usepackage[margin=1in]{geometry}
\usepackage{amsmath,amssymb,amsthm}
\usepackage{mathtools}
\usepackage{booktabs}
\usepackage{enumitem}
\usepackage{listings}
\usepackage{xcolor}
\usepackage{hyperref}
\usepackage{titlesec}
\usepackage{fancyvrb}
\usepackage{microtype}
\usepackage{parskip}
\usepackage[font=small,labelfont=bf]{caption}

% --- Hyperref setup ---
\hypersetup{
  colorlinks=true,
  linkcolor=blue!50!black,
  urlcolor=blue!50!black,
  citecolor=blue!50!black,
  pdftitle={A Proof Brokerage Architecture: Specification v1.0},
  pdfauthor={Specification working draft}
}

% --- Listings setup for code/JSON/IR ---
\definecolor{codebg}{rgb}{0.97,0.97,0.97}
\definecolor{codekw}{rgb}{0.0,0.0,0.55}
\definecolor{codestr}{rgb}{0.55,0.0,0.0}
\definecolor{codecomment}{rgb}{0.25,0.45,0.25}

\lstdefinelanguage{json}{
  basicstyle=\ttfamily\footnotesize,
  numbers=none,
  showstringspaces=false,
  breaklines=true,
  frame=single,
  framesep=4pt,
  rulecolor=\color{black!20},
  backgroundcolor=\color{codebg},
  literate=
   *{0}{{{\color{codestr}0}}}{1}
    {1}{{{\color{codestr}1}}}{1}
    {2}{{{\color{codestr}2}}}{1}
    {3}{{{\color{codestr}3}}}{1}
    {4}{{{\color{codestr}4}}}{1}
    {5}{{{\color{codestr}5}}}{1}
    {6}{{{\color{codestr}6}}}{1}
    {7}{{{\color{codestr}7}}}{1}
    {8}{{{\color{codestr}8}}}{1}
    {9}{{{\color{codestr}9}}}{1}
    {:}{{{\color{codekw}{:}}}}{1}
    {,}{{{\color{codekw}{,}}}}{1}
    {\{}{{{\color{codekw}{\{}}}}{1}
    {\}}{{{\color{codekw}{\}}}}}{1}
    {[}{{{\color{codekw}{[}}}}{1}
    {]}{{{\color{codekw}{]}}}}{1},
}

\lstdefinestyle{lean}{
  basicstyle=\ttfamily\footnotesize,
  numbers=none,
  showstringspaces=false,
  breaklines=true,
  frame=single,
  framesep=4pt,
  rulecolor=\color{black!20},
  backgroundcolor=\color{codebg},
  keywordstyle=\color{codekw}\bfseries,
  commentstyle=\color{codecomment}\itshape,
  stringstyle=\color{codestr},
  morekeywords={example,theorem,lemma,def,by,fun,if,then,else,let,in,
                 match,with,variable,instance,class,structure,inductive,
                 abbrev,private,axiom,opaque,partial,unsafe,
                 Type,Prop,Sort,fixes,assumes,shows,proof,qed},
  morecomment=[l]{--},
  morecomment=[s]{/-}{-/},
  morestring=[b]",
}

\lstdefinestyle{plain}{
  basicstyle=\ttfamily\footnotesize,
  numbers=none,
  showstringspaces=false,
  breaklines=true,
  frame=single,
  framesep=4pt,
  rulecolor=\color{black!20},
  backgroundcolor=\color{codebg},
}

% --- Theorem-like environments ---
\theoremstyle{definition}
\newtheorem{definition}{Definition}[section]
\newtheorem{example}{Example}[section]
\newtheorem{requirement}{Requirement}[section]

\theoremstyle{remark}
\newtheorem*{remark}{Remark}
\newtheorem*{rationale}{Rationale}

% --- Section formatting ---
\titleformat{\section}[hang]
  {\normalfont\Large\bfseries}
  {\thesection}{0.6em}{}
\titleformat{\subsection}[hang]
  {\normalfont\large\bfseries}
  {\thesubsection}{0.6em}{}

% --- Title ---
\title{
  \vspace{-1cm}
  \textbf{A Proof Brokerage Architecture}\\[0.3em]
  \large{Specification v1.0 (working draft)}\\[0.3em]
  \normalsize{An interchange and dispatch layer for proof assistants and automated reasoning backends}
}
\author{}
\date{Working draft — \today}

\begin{document}
\maketitle
\thispagestyle{empty}

% Global sloppiness for breaking long \texttt{...} runs
\sloppy
\emergencystretch=3em

\begin{abstract}
\noindent
This document specifies an architecture for a system that brokers proof obligations
between an interactive proof assistant (the \emph{home system}) and a
heterogeneous collection of automated reasoning \emph{backends}, including SMT
solvers, automated theorem provers (ATPs), other interactive theorem provers
(ITPs), and large-language-model (LLM) provers. The system is structured as a
three-stage pipeline: serialization of a proof obligation into a backend-independent
intermediate representation (IR); a configurable rewriting phase that performs
target-independent IR transformations; and per-backend adapters that refine the
rewritten IR into backend-native input. Returned proofs are encoded as
tiered \emph{certificates} whose verification produces home-system proof terms
through a combination of native witness checking, hint-based reconstruction,
trace replay, and (optionally) LLM-assisted translation. The architecture
deliberately preserves the home system's kernel as the trusted base for all
tiers except an explicit oracle tier; LLMs and other proposers are accepted only
through kernel-checked outputs. This document specifies the IR, the certificate
format, the dispatch protocol, the verifier, and the reconstruction layer in
sufficient detail to support a v1 implementation targeting Lean~4 as primary
home system, with Rocq (Coq) as an architectural probe.
\end{abstract}

\tableofcontents
\newpage

\section{Introduction and motivation}
\label{sec:intro}

\subsection{Problem statement}

Modern interactive proof assistants (Lean~4, Rocq, Isabelle, Agda, and others)
each ship with their own foundational kernel, library ecosystem, tactic
language, and automation tooling. A proof of a given theorem in one system
is generally not transferable to another: foundational differences (dependent
type theory vs.\ classical higher-order logic), library divergence
(\texttt{mathlib} vs.\ Isabelle's AFP vs.\ \texttt{mathcomp}), and tactic-language
incompatibilities make cross-system interchange difficult in principle and
prohibitive in practice.

Separately, the landscape of automated reasoning tools available to proof
assistant users is uneven. SMT solvers (Z3, CVC5) are powerful for arithmetic,
bitvector, and array reasoning but require careful integration to be sound.
Automated theorem provers (Vampire, E, SPASS) are excellent for first-order
logic with equality. LLM-based proof generation tools have become viable
proposers for goals where traditional automation struggles. Each of these
classes has different input formats, different output formats, different
notions of what constitutes a proof, and different relationships to the
foundational logic of the host proof assistant.

The status quo is that each home system maintains its own tactical bridges to
each backend it cares about. Lean has its own SMT integration; Rocq has
SMTCoq; Isabelle has \texttt{sledgehammer}. These bridges are excellent within
their scope but they are not interoperable: a bridge built for Isabelle does
not directly help a Lean user, and vice versa. The engineering effort to
maintain $N \times M$ bridges between $N$ home systems and $M$ backends grows
quadratically.

\subsection{Goals}

This document specifies an architecture intended to:

\begin{enumerate}[leftmargin=2em]
  \item Provide a single interchange format and dispatch pipeline through
        which proof obligations from any participating home system can be
        routed to any participating backend.
  \item Preserve the trusted base of each home system: backends propose,
        the home system's kernel disposes. Soundness is never delegated.
  \item Acknowledge and exploit the heterogeneity of backends rather than
        forcing them into a single mold. Different backends consume different
        translations of the same goal; the architecture accommodates this.
  \item Accommodate LLM-based provers as first-class backends without
        compromising soundness.
  \item Surface dispatch behavior, success rates, and failure modes to the
        user with sufficient granularity that fluent use of the system is
        learnable.
\end{enumerate}

\subsection{Non-goals}

This document does \emph{not} specify:

\begin{itemize}[leftmargin=2em]
  \item A universal foundational logic into which all home systems embed.
        Prior work on logical frameworks (notably Dedukti) explores this
        direction; the present architecture is orthogonal and complementary.
  \item Library-level translation between home systems (porting
        \texttt{mathlib} to Rocq, etc.). This is treated as a separate problem.
  \item Automatic translation of arbitrary tactic scripts between home
        systems. The unit of dispatch is a subgoal, not a tactic script.
\end{itemize}

\subsection{Architectural framing}

The system is conceived as a \emph{proof brokerage}: a market of provers
competing to discharge obligations issued by a home system. The metaphor is
deliberate. The user retains agency over which backends they employ and on
what terms. The dispatcher operates as an intermediary, not a black box.
Failed dispatches are first-class events with structured diagnostics, and
backend performance is observable through a session-level dashboard.

The architecture rejects two simpler designs. First, a single
backend-independent IR consumed uniformly by all backends (the LLVM analogy)
fails because backends have genuinely incompatible input requirements: a
goal involving a typeclass-quantified type cannot be simultaneously
specialized to linear arithmetic for an SMT solver and axiomatized
abstractly for an ATP. Second, a per-(home-system, backend) translator
matrix fails on engineering grounds: it grows quadratically and produces no
reusable infrastructure.

The chosen design is a three-stage pipeline:

\begin{enumerate}[leftmargin=2em]
  \item \textbf{Serialization.} The home system produces a single,
        backend-independent IR document that records the goal's structure
        plus enough metadata for adapters to perform target-specific
        refinements.
  \item \textbf{IR rewriting.} A configurable pipeline of target-independent
        IR-to-IR transformations simplifies the document. Quotient elimination,
        definition unfolding, and propositional simplification are examples.
        The pipeline is configured per dispatch, with adapters declaring
        their preferred preprocessing.
  \item \textbf{Adapter refinement.} A backend-specific adapter consumes
        the rewritten IR and produces the backend's native input. The
        adapter records its refinement choices in a \emph{refinement record}
        attached to any returned certificate.
\end{enumerate}

This is structurally analogous to a modern compiler pipeline (source $\to$
IR $\to$ target-independent passes $\to$ target-dependent lowering $\to$
machine code). The key adaptation is that \emph{certificates flow back}:
the proof returned by the backend must be lifted, through the inverse of
adapter refinement and IR rewriting, into a home-system proof term.

\subsection{Document structure}

Section~\ref{sec:scope} fixes the scope and milestones for v1. Section~\ref{sec:overview}
gives the architectural overview with diagrams. Section~\ref{sec:ir} specifies the
intermediate representation in detail. Section~\ref{sec:rewriting} specifies the
IR rewriting layer. Section~\ref{sec:certificates} specifies the certificate format
and the tiered verification semantics. Section~\ref{sec:dispatch} specifies the
dispatch protocol. Section~\ref{sec:verifier} specifies the certificate verifier.
Section~\ref{sec:reconstruction} specifies the reconstruction layer.
Section~\ref{sec:llm} discusses LLM integration points. Section~\ref{sec:itp-itp}
discusses ITP-to-ITP dispatch and its restrictions in v1. Section~\ref{sec:examples}
walks through three worked examples end-to-end. Appendix~\ref{app:registry}
specifies the shared registry of recognized patterns. Appendix~\ref{app:open}
lists open questions deferred to v2 or beyond.

\section{Scope and milestones for v1}
\label{sec:scope}

\subsection{Home systems}

\begin{itemize}[leftmargin=2em]
  \item \textbf{Lean~4 (primary).} Full integration: serializer, all enabled
        verifier tiers, all reconstruction strategies, build-time integration,
        and the session dashboard.
  \item \textbf{Rocq (architectural probe).} IR serializer and Tier~0 plus
        Tier~1 verifiers only. No Tier~2 or Tier~3 reconstruction, no LLM
        strategies, no build-time integration. The Rocq integration exists
        to validate that the IR and protocol are not Lean-shaped, and to
        round-trip simple obligations. Production-quality Rocq integration
        is a v2 milestone.
\end{itemize}

\subsection{Backend categories}

\begin{itemize}[leftmargin=2em]
  \item \textbf{SMT.} Initial target: CVC5 (with Alethe proof output). The
        adapter exposes hooks for Z3 and other SMT-LIB-2 conformant solvers.
  \item \textbf{ATP.} Initial target: Vampire (with TPTP proof output).
        The adapter exposes hooks for E and other TPTP-conformant provers.
  \item \textbf{LLM provers.} A generic LLM-prover adapter that accepts a
        configured model endpoint and produces tactic scripts in the home
        system's language. Examples: DeepSeek-Prover-class systems,
        general-purpose LLMs prompted appropriately.
  \item \textbf{ITPs (restricted).} ITP-to-ITP dispatch is supported with
        Tier~0 (oracle, with explicit trust opt-in) or Tier~2 (lemma list)
        only. See Section~\ref{sec:itp-itp}.
\end{itemize}

\subsection{Tiers supported in v1}

\begin{itemize}[leftmargin=2em]
  \item \textbf{Tier~0 (oracle).} Supported, gated behind explicit
        per-project trust configuration.
  \item \textbf{Tier~1 (witness).} Supported for: linear arithmetic
        (Farkas certificates), propositional satisfiability (assignments
        and unsat cores). Other witness kinds deferred.
  \item \textbf{Tier~2 (lemma list).} Supported as a \emph{fallback tier};
        framing and demotion clause described in Section~\ref{sec:certificates}.
  \item \textbf{Tier~3 (proof trace).} Supported for Alethe (CVC5) and
        TPTP-FOF/THF (Vampire). Strategy stack includes symbolic replay,
        cross-system tactic translation, and optional LLM-assisted
        reconstruction. Other formats deferred.
  \item \textbf{Tier~4 (foundational certificate).} Reserved; not
        implemented in v1.
\end{itemize}

\subsection{Out of scope for v1}

\begin{itemize}[leftmargin=2em]
  \item Production-quality Rocq integration (deferred to v2).
  \item Tier~3 reconstruction for backend formats other than Alethe and
        TPTP.
  \item Tier~4 foundational certificates and the associated logical
        framework embedding work.
  \item ML-driven premise selection (symbol-based selection only in v1).
  \item Speculative dispatch with predictive caching across user edits.
  \item Library-level translation and porting.
  \item Conversion metadata (definitional equality conventions, normalization
        regime). The v1 IR assumes terms are presented in a fully
        $\beta$-$\eta$-normalized form; goals depending on non-normalized
        distinctions are not supported in v1.
\end{itemize}

\section{Architectural overview}
\label{sec:overview}

\subsection{Components}

The system comprises seven major components, illustrated in
Figure~\ref{fig:overview}.

\begin{figure}[h]
\centering
\small
\begin{verbatim}
   +-----------------+
   |   Home System   |  Lean 4 (primary) | Rocq (probe)
   |     Plugin      |  - serializer
   |                 |  - reconstruction layer
   +--------+--------+  - dashboard
            |
            v
   +-----------------+
   |   IR Document   |  general, backend-independent
   +--------+--------+
            |
            v
   +-----------------+
   |   IR Rewriter   |  configurable per-dispatch pipeline
   |                 |  - quotient elimination, unfolding,
   +--------+--------+    propositional simplification, ...
            |
            v
   +-----------------+
   |   Dispatcher    |  capability matching, backend selection
   +--------+--------+
            |
            v
   +-----------------+
   | Backend Adapter |  per-backend refinement, native input,
   |                 |  refinement record generation
   +--------+--------+
            |
            v
   +-----------------+      +-----------------+
   |  Backend (CVC5, |----->|  Certificate    |
   |  Vampire, LLM,  |      |     (with       |
   |   Isabelle)     |      |  refinement     |
   +-----------------+      |    record)      |
                            +--------+--------+
                                     |
                                     v
                            +-----------------+
                            |   Certificate   |
                            |    Verifier     |
                            +--------+--------+
                                     |
                                     v
                            +-----------------+
                            | Reconstruction  |
                            |     Layer       |  (back to home system)
                            +-----------------+
\end{verbatim}
\caption{High-level architectural overview. The forward pipeline (top to
bottom) takes a home-system proof state to a backend-native input. The
return pipeline (bottom-right back to the home system) verifies the
returned certificate and reconstructs a home-system proof term.}
\label{fig:overview}
\end{figure}

\begin{description}
  \item[Home System Plugin.] A tactic and metaprogram extension installed
        in the home system. Exposes the \texttt{dispatch} family of tactics,
        serializes proof state into IR, and consumes verifier outputs to
        produce home-system proof terms. Per-home-system; the Lean and Rocq
        plugins are distinct artifacts but share the IR schema.
  \item[IR Rewriter.] A configurable pipeline of IR-to-IR passes operating
        on the general IR. Each pass is target-independent; the choice of
        passes is configured per dispatch, typically driven by adapter
        preferences. Implemented as a shared library callable from any
        component.
  \item[Dispatcher.] Decides which backend(s) to invoke based on goal
        classification, backend capabilities, user directives, and
        historical performance. Coordinates parallel speculative dispatch
        when configured.
  \item[Backend Adapter.] Per-backend translator from rewritten IR to
        backend-native input. Generates a refinement record describing its
        translation choices. Receives the backend's output and packages it
        as a certificate.
  \item[Certificate Verifier.] Tier-dispatched verification of returned
        certificates. Produces home-system proof terms or structured
        rejection reasons.
  \item[Reconstruction Layer.] Lifts verified certificates through the
        inverse of adapter refinement and IR rewriting to produce a proof
        term of the original home-system goal type.
  \item[Certificate Cache.] Project-local store keyed on goal hash plus
        dispatch context plus backend identity. Invalidated on backend
        version change or referenced-lemma content hash change.
\end{description}

\subsection{Soundness story}

The trusted base for Tiers 1, 2, and 3 is the home system's kernel.
Witnesses, reconstruction outputs, and replayed traces are all packaged as
home-system proof terms before being accepted; the kernel re-checks every
such term. No backend, LLM, or external service is in the trusted base for
these tiers.

Tier~0 (oracle) is the sole exception: it expands the trusted base to
include the named backend's kernel, and is gated behind explicit per-project
trust configuration. Tier~0 is intended for cases where the user has made
an informed trust decision, such as accepting an Isabelle-proved lemma
during cross-ITP dispatch.

LLMs participate only as proposers. An LLM-generated tactic script is
accepted iff the home system's kernel accepts it; an LLM-translated proof
trace is likewise kernel-checked. The architecture provides no mechanism
by which LLM output bypasses kernel verification.

\subsection{Two paths through the system}

\paragraph{Slow path (interactive).} The user invokes a \texttt{dispatch}
tactic. The plugin serializes the proof state, runs the IR rewriter, hands
off to the dispatcher, awaits a result, and (on success) inserts the
reconstructed proof term. UI feedback is provided asynchronously: the
tactic call shows a pending indicator until the dispatch resolves, with
per-backend timing visible in the dashboard. Typical latency: 1--30
seconds.

\paragraph{Build path.} On project compilation, dispatch calls re-run with
longer timeout budgets and persistent caching. The build-time path is the
source of truth for CI: a project's CI passes iff every dispatch call
produces a kernel-checked proof term within budget. Cached certificates
from prior interactive runs are consulted before re-dispatch.

\section{The Intermediate Representation}
\label{sec:ir}

\subsection{Design principles}

The IR is the document produced by the home-system serializer and consumed
by the IR rewriter and (after rewriting) by backend adapters. It is
designed to satisfy the following principles:

\begin{description}
  \item[Backend-independence.] The IR is produced once, before backend
        selection. The same IR document is, in principle, consumable by
        any compatible backend.
  \item[Bounded metadata categories.] The IR's metadata is organized into
        four categories (logic classification, type metadata, definitional
        metadata, library provenance). New language features extend the
        \emph{content} of these categories but do not introduce new
        categories.
  \item[Self-describing.] Every IR document declares its schema version,
        its source system, and its context tier. Consumers can
        validate compatibility statically.
  \item[Round-trippable.] Every translation choice made by the serializer
        is recorded with sufficient information that reconstruction can
        invert it. The IR is not lossy with respect to the home-system
        proof state.
  \item[No source-system-specific structure.] The IR records abstract
        signatures and shape descriptors, not Lean-specific instance
        resolution traces or Rocq-specific canonical structure data.
\end{description}

\subsection{Top-level structure}

An IR document is a JSON value with the following top-level fields. (Binary
serialization via CBOR is the recommended production format; JSON is the
canonical specification format.)

\begin{lstlisting}[language=json]
{
  "ir_version": "1.0",
  "source_system": { "name": "lean", "version": "4.x.y" },
  "tier": "structural",
  "logic_classification": <LogicClassification>,
  "goal": <Goal>,
  "context": <Context>,
  "type_metadata": { "<type-id>": <TypeMetadata>, ... },
  "definitional_metadata": { "<symbol-id>": <DefinitionalMetadata>, ... },
  "library_provenance": { "<entity-name>": <Provenance>, ... },
  "user_directives": <UserDirectives>
}
\end{lstlisting}

The \texttt{tier} field takes values \texttt{"goal"}, \texttt{"structural"},
or \texttt{"rich"}, controlling how much context is included
(see Section~\ref{sec:ir-tiers}).

\subsection{Logic classification}
\label{sec:ir-logic}

The \texttt{logic\_classification} field tells dispatchers and adapters
what kind of logical content the goal contains, so that
incompatible backends can be filtered out before invocation.

\begin{lstlisting}[language=json]
{
  "order": "first_order" | "higher_order" | "dependent",
  "features_used": [ <feature-tag>, ... ],
  "first_order_fragment": <fragment-id> | "none",
  "decidable_theory": <theory-id> | null
}
\end{lstlisting}

The \texttt{order} field is the highest order of quantification appearing
in the goal or its hypotheses. The \texttt{features\_used} list draws on
the registry (Appendix~\ref{app:registry}) of recognized logical features:
\texttt{function\_quantification}, \texttt{predicate\_argument},
\texttt{equality\_in\_quotient\_type}, \texttt{induction\_recursion}, and
others. The \texttt{first\_order\_fragment} field, when present, names a
decidable or semi-decidable fragment that contains the goal: \texttt{LIA},
\texttt{LRA}, \texttt{NIA}, \texttt{BV}, \texttt{UF}, \texttt{UFLIA}, etc.,
or \texttt{"none"} if no such fragment applies. The
\texttt{decidable\_theory} field, when set, asserts that the goal lies in
a particular decidable theory; backends may use this to dispatch directly
to a decision procedure.

\begin{requirement}[Logic classification soundness]
The serializer must not assign a \texttt{first\_order\_fragment} narrower
than the actual fragment of the goal. False classification of a
higher-order goal as first-order is a soundness-relevant bug; backends
relying on the classification may produce certificates that fail to lift
back. Conservative over-approximation (claiming a wider fragment than
necessary) is acceptable but may reduce dispatch effectiveness.
\end{requirement}

\subsection{The goal and the shell calculus}
\label{sec:ir-shell}

The \texttt{goal} field encodes the proposition to be proved as a
\emph{shell term} together with optional \emph{payloads}.

\begin{lstlisting}[language=json]
{
  "shell": <ShellTerm>,
  "payloads": { "<payload-id>": <SourceTerm>, ... }
}
\end{lstlisting}

The shell calculus is a logic-agnostic typed term calculus designed to
admit faithful translation from any participating home system's
proposition language. It is not a foundational logic; it is a structural
representation. The grammar of shell terms is:

\begin{align*}
\textit{ShellTerm} ::= \;\;& \textsf{Forall}(\textit{var}, \textit{type}, \textit{ShellTerm}) \\
                       \mid \;& \textsf{Exists}(\textit{var}, \textit{type}, \textit{ShellTerm}) \\
                       \mid \;& \textsf{Lambda}(\textit{binders}, \textit{ShellTerm}) \\
                       \mid \;& \textsf{Implies}(\textit{ShellTerm}, \textit{ShellTerm}) \\
                       \mid \;& \textsf{And}(\textit{ShellTerm}, \textit{ShellTerm}) \\
                       \mid \;& \textsf{Or}(\textit{ShellTerm}, \textit{ShellTerm}) \\
                       \mid \;& \textsf{Not}(\textit{ShellTerm}) \\
                       \mid \;& \textsf{Eq}(\textit{type}, \textit{ShellTerm}, \textit{ShellTerm}) \\
                       \mid \;& \textsf{App}(\textit{symbol}, \textit{type-args}, \textit{args}) \\
                       \mid \;& \textsf{Var}(\textit{name}) \\
                       \mid \;& \textsf{Const}(\textit{name}) \\
                       \mid \;& \textsf{NumLit}(\textit{value}) \\
                       \mid \;& \textsf{Opaque}(\textit{payload-id})
\end{align*}

The \textsf{Eq} constructor carries the type at which equality is
asserted, because equality in dependent type theory is type-indexed.
The \textsf{App} constructor takes type arguments separately from term
arguments, supporting polymorphic and higher-kinded application without
encoding type abstraction in the shell.

The \textsf{Opaque} constructor is the escape hatch. Any source-language
construct that cannot be expressed in the shell (inductive type
eliminators with complex motives, dependent pattern matches, etc.) is
encapsulated as an opaque payload with a fresh identifier; the original
source term is recorded under that identifier in the \texttt{payloads}
field. Backends that understand the source system's term language may
unpack opaque payloads; backends that do not treat them as
uninterpreted constants of an inferred type.

\begin{remark}
The shell deliberately does \emph{not} encode bound-variable conventions
(de~Bruijn indices vs.\ named variables), eta-conversion behavior, or
universe levels in any but the most superficial way. These are
home-system concerns that the IR avoids committing to. Where
universe-polymorphism information is needed (rare in v1), it is recorded
in type metadata, not in the shell.
\end{remark}

\subsection{Context and tiers}
\label{sec:ir-tiers}

The \texttt{context} field encodes hypotheses and free variables. Three
context tiers control the information included.

\paragraph{Tier ``goal''.} Only the proposition. No hypotheses, no free
variables. Used when the goal is self-contained.

\paragraph{Tier ``structural''.} Hypotheses translated as shell terms,
free variables typed via references to type metadata, but no library
lemmas. Default for most ATP and SMT dispatches.

\paragraph{Tier ``rich''.} Structural context plus a curated library
slice. Lemmas in the slice are themselves shell terms with provenance.
For v1, the slice is constructed by:

\begin{itemize}[leftmargin=2em]
  \item \textbf{Symbol-based premise selection.} Lemmas whose
        statements share head symbols (or close generalizations) with
        the goal are included automatically.
  \item \textbf{Manual augmentation.} Lemmas explicitly cited via the
        \texttt{using} clause of the dispatch tactic are included
        unconditionally.
  \item \textbf{No transitive closure.} Only lemma statements are
        included; their proofs and dependencies are not recursively
        translated.
  \item \textbf{Dropped lemmas are reported.} Lemmas whose statements
        cannot be translated under the structural shell are omitted from
        the slice with a diagnostic, not silently retained.
\end{itemize}

ML-driven premise selection is deferred to v2.

\subsection{Type metadata}
\label{sec:ir-type-metadata}

For every type appearing in the goal or context that is not a primitive
of the shell, the \texttt{type\_metadata} field records information
sufficient for adapters to reason about that type abstractly. The
schema is:

\begin{lstlisting}[language=json]
{
  "kind": "primitive" | "type_variable" | "type_constructor_application",
  ... (kind-specific fields) ...
}
\end{lstlisting}

\subsubsection{Kind: \texttt{primitive}}

For known primitive types (\texttt{Nat}, \texttt{Int}, \texttt{Bool}, \texttt{String}, etc.):

\begin{lstlisting}[language=json]
{
  "kind": "primitive",
  "name": "Int",
  "theory_tags": ["arithmetic", "ring", "linearly_ordered"]
}
\end{lstlisting}

The \texttt{theory\_tags} draw from the registry; they are how
adapters recognize that a type is suitable for a particular backend
theory.

\subsubsection{Kind: \texttt{type\_variable}}

For type variables introduced by the proof state (e.g., a polymorphic
goal):

\begin{lstlisting}[language=json]
{
  "kind": "type_variable",
  "name": "<alpha-id>",
  "kind_class": "Type",
  "instances": [
    {
      "instance_name": "<inst-id>",
      "class": {
        "name": "LinearOrderedCommRing",
        "library": "mathlib",
        "library_version": "4.x.y",
        "content_hash": "sha256:..."
      },
      "abstract_signature": <AbstractSignature>,
      "abstract_axioms": [ <AxiomShape>, ... ],
      "theory_classification_tags": [ <tag>, ... ]
    }
  ]
}
\end{lstlisting}

The \texttt{abstract\_signature} describes the operations of the
typeclass without committing to home-system-specific resolution. The
\texttt{abstract\_axioms} use shape descriptors from the registry
(see Appendix~\ref{app:registry}, Section~\ref{app:axiom-shapes}). The
\texttt{theory\_classification\_tags} record valid embeddings into
backend theories: e.g., \texttt{embeds\_into:Int\_for\_universal\_LIA}
asserts that universal LIA sentences valid over $\mathbb{Z}$ are valid
over any instance of this class.

\begin{requirement}[Theory tag soundness]
A \texttt{theory\_classification\_tag} must correspond to a verified
embedding. The serializer is responsible for issuing only tags that the
home system can justify with an embedding lemma in its library; the
lemma reference is recorded in library provenance. Spurious tags are
soundness-relevant bugs.
\end{requirement}

\subsubsection{Kind: \texttt{type\_constructor\_application}}

For types built by a non-primitive type former (quotient, inductive,
sigma, refinement, higher-inductive):

\begin{lstlisting}[language=json]
{
  "kind": "type_constructor_application",
  "constructor": {
    "name": "<constructor-id>",
    "construction_kind":
      "quotient" | "inductive" | "sigma" | "refinement" | "higher_inductive",
    ... (construction-kind-specific fields) ...
  },
  "arguments": [ <ShellTerm>, ... ]
}
\end{lstlisting}

For \texttt{construction\_kind: "quotient"}, the constructor block
includes \texttt{underlying\_type}, \texttt{equivalence\_relation}
(with shell rendering and proof of equivalence), an
\texttt{elimination\_principle} (the induction rule), and an
\texttt{equality\_principle} (the soundness rule, e.g., \texttt{Quot.sound}).

For \texttt{construction\_kind: "inductive"}, the constructor block
includes a list of \texttt{constructors} (each with a name, type
schema, and recursive position annotations) and an
\texttt{elimination\_principle} describing the recursor.

Other construction kinds follow analogous patterns.

\subsection{Definitional metadata}
\label{sec:ir-defn-metadata}

For every defined symbol appearing in the goal or context that is not a
shell primitive, the \texttt{definitional\_metadata} field records
information sufficient for adapters to reason about that symbol. The
schema is:

\begin{lstlisting}[language=json]
{
  "kind":
    "primitive_arithmetic" | "typeclass_method" | "defined_function" |
    "lifted_to_quotient" | "constructor" | "eliminator",
  ... (kind-specific fields) ...
}
\end{lstlisting}

\subsubsection{Kind: \texttt{primitive\_arithmetic}}

\begin{lstlisting}[language=json]
{
  "kind": "primitive_arithmetic",
  "abstract_signature": "Int -> Int -> Int",
  "theory_tag": "LIA:plus"
}
\end{lstlisting}

\subsubsection{Kind: \texttt{typeclass\_method}}

\begin{lstlisting}[language=json]
{
  "kind": "typeclass_method",
  "method_name": "HAdd.hAdd",
  "host_class": "<class-id>",
  "abstract_role": "addition" | "ordering" | "literal_inclusion" | ...,
  "specialization_targets": [
    { "theory": "LIA", "operator": "+" },
    { "theory": "LRA", "operator": "+" }
  ]
}
\end{lstlisting}

\subsubsection{Kind: \texttt{defined\_function}}

\begin{lstlisting}[language=json]
{
  "kind": "defined_function",
  "abstract_signature": "(B -> C) -> (A -> B) -> (A -> C)",
  "definitional_equation": <ShellTerm>,
  "extensional_axiom": <ShellTerm>,
  "concept_tag": "function_composition"
}
\end{lstlisting}

The \texttt{definitional\_equation} expresses the function's defining
rewrite in shell form, suitable for unfolding by the IR rewriter or
adapter. The \texttt{concept\_tag}, when present, ties the symbol to a
concept in the registry, allowing adapters to recognize standard
mathematical operations across home systems.

\subsubsection{Kind: \texttt{lifted\_to\_quotient}}

\begin{lstlisting}[language=json]
{
  "kind": "lifted_to_quotient",
  "host_type": "<type-id>",
  "underlying_function": {
    "name": "Int.add",
    "shell": <ShellTerm>
  },
  "lifting_obligation": {
    "shape": "respects_relation",
    "discharged_at": "definition_site"
  }
}
\end{lstlisting}

The IR rewriter's \emph{quotient elimination} pass uses this metadata to
replace lifted-function applications with applications of the underlying
function under quotient introduction, when this transformation is
beneficial for the target backend.

\subsubsection{Other kinds}

\texttt{constructor} and \texttt{eliminator} kinds record information for
inductive type constructors and recursors respectively. Their schemas
parallel \texttt{lifted\_to\_quotient} and are described in
Appendix~\ref{app:registry}.

\subsection{Library provenance}

\begin{lstlisting}[language=json]
{
  "<entity-name>": {
    "library": "mathlib",
    "version": "4.x.y",
    "module_path": "Mathlib.Algebra.Order.Ring.Defs",
    "content_hash": "sha256:..."
  }
}
\end{lstlisting}

Provenance is recorded for every named entity (typeclass, defined
function, lemma) referenced from the IR. The reconstruction layer uses
the content hash to validate that the home system has the expected
version of each entity available before consuming a certificate; a hash
mismatch invalidates the certificate.

\subsection{User directives}

\begin{lstlisting}[language=json]
{
  "preferred_backend": "cvc5" | "vampire" | "llm" | null,
  "tier_preference": ["1", "3", "2", "0"],
  "rewriter_preferences": {
    "enable_quotient_elimination": true,
    "enable_definition_unfolding": ["function_composition", ...],
    "disable_passes": [...]
  },
  "budget": { "wall_time_ms": 30000, "memory_mb": 4096 }
}
\end{lstlisting}

User directives capture preferences supplied through the dispatch
tactic's named arguments (Section~\ref{sec:dispatch}).

\subsection{Translation discipline}
\label{sec:ir-discipline}

The home-system serializer is bound by the following discipline:

\begin{requirement}[No unjustified specialization]
The serializer must not assign a \texttt{theory\_classification\_tag} or
generate a \texttt{specialization\_targets} entry without a corresponding
verified embedding lemma in the home system's library, recorded in
provenance.
\end{requirement}

\begin{requirement}[Conservative classification]
When in doubt, the serializer assigns a wider logic classification, a
broader fragment, or fewer theory tags. False positives in
classification are soundness-relevant; false negatives merely reduce
effectiveness.
\end{requirement}

\begin{requirement}[Round-trip information sufficiency]
For any IR document the serializer produces, the metadata must be
sufficient that any home-system proof term claiming the document's
\texttt{goal} (in shell form) can be lifted by the reconstruction layer
to a home-system proof term of the original proof state's goal type.
\end{requirement}

This last requirement is the core soundness obligation of the
home-system plugin and is the contract under which all other
architecture components operate.

\section{The IR rewriting layer}
\label{sec:rewriting}

\subsection{Purpose and architecture}

The IR rewriting layer performs target-independent transformations on
IR documents. Its purpose is to factor out simplifications that are
useful across multiple backends, so that adapters do not redundantly
implement the same logic. Examples of such transformations:

\begin{itemize}[leftmargin=2em]
  \item \textbf{Quotient elimination.} A goal stated as an equality in a
        quotient type is reduced to a goal at the underlying type using
        the elimination and equality principles, when the metadata
        permits.
  \item \textbf{Definition unfolding.} Specific defined functions are
        replaced with their definitional equations, exposing the
        underlying mathematical content. Unfolding is selective: not all
        definitions should be unfolded, and some are unfolded only for
        certain backends.
  \item \textbf{Propositional simplification.} Boolean simplifications
        ($P \wedge \top \leadsto P$, $P \vee \neg P \leadsto \top$, etc.)
        and trivial restructuring.
  \item \textbf{Equality saturation.} Bounded equality reasoning to
        canonicalize terms before backend dispatch.
  \item \textbf{Quantifier prenexing.} Moving quantifiers outward,
        useful for backends that prefer prenex normal form.
  \item \textbf{Skolemization.} Eliminating existentials by introducing
        fresh function symbols.
\end{itemize}

Each rewrite is implemented as a \emph{pass}: a function from IR
documents to IR documents that preserves logical equivalence (modulo
the soundness conditions documented for each pass). Passes are
composed into pipelines.

\subsection{Pass interface}

A pass is specified by:

\begin{lstlisting}[language=json]
{
  "name": "<pass-name>",
  "version": "1.0",
  "preconditions": [ <precondition>, ... ],
  "postconditions": [ <postcondition>, ... ],
  "soundness_obligation": "<text-or-reference>",
  "configuration_schema": <ConfigSchema>
}
\end{lstlisting}

Preconditions name properties an IR document must satisfy for the pass
to be applicable; postconditions name properties the output is
guaranteed to satisfy. The soundness obligation states what makes the
pass logically valid (typically a one-line statement of the form ``the
output IR is provable iff the input IR is provable, given certain
metadata is correct'').

\begin{requirement}[Pass soundness preservation]
Every pass must preserve provability of the goal. If the input IR
encodes a goal $G$ in context $\Gamma$, the output IR must encode a goal
$G'$ in context $\Gamma'$ such that $\Gamma \vdash G$ if and only if
$\Gamma' \vdash G'$, in the home system's foundational logic.
\end{requirement}

\subsection{Configurability}

The rewriting pipeline is configured per dispatch. Configuration is
determined by, in priority order: (a) explicit user directives in the
IR document; (b) adapter preferences declared by the dispatching
backend; (c) project-level defaults; (d) global defaults.

\begin{lstlisting}[language=json]
{
  "pipeline": [
    { "pass": "definition_unfolding",
      "config": { "concepts": ["function_composition"] } },
    { "pass": "quotient_elimination",
      "config": { "aggressive": false } },
    { "pass": "propositional_simplification" }
  ],
  "stop_on_failure": false,
  "timeout_per_pass_ms": 1000
}
\end{lstlisting}

\subsection{The default pipeline}

The default rewriting pipeline for v1, applied when no specific
configuration is provided:

\begin{enumerate}[leftmargin=2em]
  \item \texttt{propositional\_simplification}: minimal, removes trivial
        constructs.
  \item \texttt{definition\_unfolding}: only for symbols whose
        \texttt{definitional\_metadata} declares \texttt{concept\_tag} from
        a registered ``always-unfold-for-dispatch'' list (function
        composition, identity, etc.).
  \item No quotient elimination by default; adapters that benefit from it
        opt in via their preferences.
  \item No prenexing or Skolemization by default; SMT and ATP adapters
        opt in.
\end{enumerate}

\subsection{Backend adapter preferences}

Each backend adapter declares a preferred rewriting pipeline that is
applied before the adapter consumes the IR. Examples:

\begin{itemize}[leftmargin=2em]
  \item \textbf{CVC5 adapter.} Prefers quotient elimination (when
        applicable), aggressive definition unfolding for arithmetic
        operations, propositional simplification, and Skolemization.
  \item \textbf{Vampire adapter.} Prefers prenex normal form,
        Skolemization, and selective definition unfolding (preserving
        symbols that benefit from FOL axiomatization rather than
        unfolding).
  \item \textbf{LLM adapter.} Prefers minimal rewriting; the LLM
        consumes goals close to source syntax.
  \item \textbf{ITP adapter (e.g., Isabelle).} Prefers minimal rewriting;
        the target ITP performs its own normalization.
\end{itemize}

When parallel speculative dispatch is configured (multiple backends in
parallel), each backend receives an IR rewritten according to its own
preferences. The cost is duplicated rewriting work; the benefit is that
each backend sees the IR in its preferred form.

\subsection{Recording rewrites for reconstruction}

Each pass that runs records a \emph{rewrite trace entry} attached to the
IR document. The trace is consulted during reconstruction to invert the
rewriting and lift backend proofs back to the original goal. A trace
entry has the form:

\begin{lstlisting}[language=json]
{
  "pass": "quotient_elimination",
  "version": "1.0",
  "before_hash": "sha256:...",
  "after_hash": "sha256:...",
  "inversion_data": <pass-specific>
}
\end{lstlisting}

The \texttt{inversion\_data} field carries information sufficient to
construct the lifting from the rewritten goal back to the original. For
quotient elimination, this includes the introduction points of the
elimination principle and the witnesses to the equivalence-relation
respect conditions; the reconstruction layer uses these to wrap a
proof of the underlying-type goal in the appropriate
\texttt{Quot.ind}/\texttt{Quot.sound} structure.

\begin{requirement}[Trace completeness]
For every rewrite pass that runs, sufficient inversion data must be
recorded that the pass can be inverted by the reconstruction layer
without re-consulting the source proof state. The reconstruction layer
operates against the rewrite trace and the certificate, not against the
home system's elaborator.
\end{requirement}

\subsection{Failure modes}

Passes can fail in two ways: (a) preconditions not satisfied, in which
case the pass is skipped; (b) internal pass logic raises an error, in
which case the pipeline either aborts or continues per configuration.
A failed pass does not invalidate previous passes' outputs; the IR
document carries the trace up to the point of failure, and the
dispatcher can either invoke the adapter on the partial output or
abandon the dispatch.

\begin{remark}
The IR rewriter is a shared library callable from any component. The
home system plugin invokes it before dispatch; backend adapters may
invoke it again with their preferences; the verifier may invoke
inversion-related queries during reconstruction. Centralizing rewriting
logic in a shared library, rather than duplicating it per adapter, is
the architectural commitment that makes Position~3 of the design
debate (general IR plus per-backend refinement) work in practice.
\end{remark}

\section{The certificate format}
\label{sec:certificates}

\subsection{Purpose}

A certificate is the artifact produced by a backend adapter on
successful dispatch. It contains the information needed to verify that
the dispatched goal was proved and to reconstruct a home-system proof
term. Certificates are organized into \emph{tiers} based on the strength
of evidence they provide.

\subsection{Envelope}

Every certificate, regardless of tier, has the following envelope:

\begin{lstlisting}[language=json]
{
  "cert_version": "1.0",
  "tier": 0 | 1 | 2 | 3 | 4,
  "format": "<format-id>",
  "goal": <Goal>,
  "dispatch_context_hash": "sha256:...",
  "rewrite_trace_hash": "sha256:...",
  "backend": {
    "name": "cvc5",
    "version": "1.2.0",
    "config_hash": "sha256:..."
  },
  "resources": {
    "wall_time_ms": 1234,
    "memory_peak_kb": 56789,
    "budget_consumed": 0.7
  },
  "refinement_record": <RefinementRecord>,
  "payload": <TierPayload>
}
\end{lstlisting}

The \texttt{goal} field carries the goal in its rewritten form (post-IR-rewriter,
post-adapter-refinement). The \texttt{dispatch\_context\_hash} and
\texttt{rewrite\_trace\_hash} let the verifier confirm that the certificate
addresses the question that was asked, and that the rewrite trace is
available for inversion.

The \texttt{refinement\_record} field captures the adapter's
translation choices, described next.

\subsection{Refinement records}
\label{sec:refinement-records}

A refinement record is a structured description of what the backend
adapter did to translate the (rewritten) IR into backend-native input.
It is what the reconstruction layer uses to lift the backend's proof
back through the adapter's translation.

\begin{lstlisting}[language=json]
{
  "adapter": "<adapter-id>",
  "adapter_version": "1.0",
  "specializations": [
    {
      "kind": "type_specialization" | "method_specialization" |
              "axiomatization" | "definition_unfolding" |
              "logic_encoding" | "rendering",
      "source": "<source-form>",
      "target": "<target-form>",
      "justification": "<reference-to-metadata-or-tag>",
      "soundness_witness": "<lemma-reference>"
    },
    ...
  ],
  "fragment": "LIA" | "LRA" | "FOL" | "HOL" | ...,
  "auxiliary": <adapter-specific>
}
\end{lstlisting}

Each \texttt{specialization} entry records one translation decision made
by the adapter. The \texttt{kind} field categorizes the decision:

\begin{description}
  \item[\texttt{type\_specialization}.] A type variable or abstract type
        was replaced with a concrete backend-supported type
        (e.g., $\alpha \to \mathbb{Z}$ via an embedding lemma).
  \item[\texttt{method\_specialization}.] A typeclass method or
        polymorphic operator was specialized to a backend-native
        operator (e.g., \texttt{HAdd.hAdd} $\to$ SMT-LIB \texttt{+}).
  \item[\texttt{axiomatization}.] A typeclass instance or defined symbol
        was encoded as an uninterpreted symbol with explicit axioms.
  \item[\texttt{definition\_unfolding}.] A defined symbol was unfolded
        using its definitional equation.
  \item[\texttt{logic\_encoding}.] The shell-level logic was encoded
        into a backend-specific logical formalism (e.g., higher-order
        shell into TPTP-THF with combinators).
  \item[\texttt{rendering}.] For LLM and ITP adapters, the goal was
        rendered as source-language syntax.
\end{description}

\begin{requirement}[Refinement record completeness]
Every adapter decision that affects the form of the goal as seen by the
backend must be recorded in the refinement record. Decisions that the
reconstruction layer cannot recover from the record are soundness-relevant:
they prevent lifting the backend's proof back to the original goal.
\end{requirement}

\subsection{Tier 0: Oracle}

\begin{lstlisting}[language=json]
"payload": {
  "claim": "proved",
  "backend_attestation": "<optional-signed-claim>"
}
\end{lstlisting}

\paragraph{Verifier behavior.} Consult the project's trust policy.
Accept iff the backend is on the project's explicit Tier~0 trust list
and the trust policy permits Tier~0 acceptance for goals matching this
shape and library context.

\paragraph{Trust expansion.} Accepting a Tier~0 certificate expands the
trusted base to include the named backend's kernel. The home-system
plugin must surface this to the user at the time of acceptance, not
silently.

\paragraph{Use cases.} Rapid prototyping, ITP-to-ITP dispatch where
no usable proof artifact exists, and cases where the user has
deliberately opted into trusting a particular backend for a particular
class of goals.

\subsection{Tier 1: Witness}

\begin{lstlisting}[language=json]
"payload": {
  "witness_kind":
    "farkas" | "sat_assignment" | "sat_unsat_core" |
    "polynomial_positivstellensatz" | ...,
  "witness_data": <kind-specific>,
  "checking_recipe": "<reference-to-home-system-checker>"
}
\end{lstlisting}

\paragraph{Verifier behavior.} Run the home-system witness checker on
the witness data. The checker is a small, audited home-system tactic
that produces a kernel-checkable proof term directly from the witness.
Success requires no LLM and no probabilistic step.

\paragraph{Trust expansion.} None. The witness checker is part of the
home system; the kernel re-checks its output.

\paragraph{Witness kinds supported in v1.}

\begin{itemize}[leftmargin=2em]
  \item \texttt{farkas}: a Farkas certificate for linear arithmetic.
        Witness data is a list of nonnegative coefficients on
        hypotheses summing to a contradiction for $\neg G$.
  \item \texttt{sat\_assignment}: a satisfying assignment for a
        propositional formula. Witness data is a variable-to-Boolean map.
  \item \texttt{sat\_unsat\_core}: a minimal unsatisfiable subset for a
        propositional formula. Witness data is a list of clause indices.
\end{itemize}

Other witness kinds (Positivstellensatz for nonlinear arithmetic, model
checking witnesses, etc.) are deferred.

\subsection{Tier 2: Lemma list / reconstruction hint (fallback)}
\label{sec:tier2}

\begin{lstlisting}[language=json]
"payload": {
  "lemmas_used": [
    { "name": "<lemma-name>",
      "library": "mathlib",
      "version": "4.x.y",
      "content_hash": "sha256:..." },
    ...
  ],
  "strategy_hint":
    "aesop" | "auto" | "simp_arith" | "smt_reconstruct" | ...,
  "structural_hint": <optional-ordering-or-shape>
}
\end{lstlisting}

\paragraph{Status: fallback tier.} Tier~2 is selected when the backend
cannot produce a structured proof trace, or when no enabled Tier~3
strategy can consume the trace.

\paragraph{Verifier behavior.} Invoke the home system's reconstruction
tactic with the listed lemmas as hints. Success iff the tactic produces
a kernel-checkable term within the reconstruction budget.

\paragraph{Soundness vs.\ completeness.}

\begin{itemize}[leftmargin=2em]
  \item \emph{Sound.} Every accepted Tier~2 certificate corresponds to a
        kernel-checked home-system proof term. Reconstruction never
        extends the trusted base.
  \item \emph{Probabilistically complete.} Reconstruction can fail even
        when the originating backend succeeded. The success-rate gap
        between ``backend found a proof'' and ``home system reconstructed
        it from the lemma list'' is real and well-documented in
        \texttt{sledgehammer}-class systems. This gap is the price of
        Tier~2's flexibility.
\end{itemize}

\paragraph{UX commitments.}

\begin{itemize}[leftmargin=2em]
  \item Successful Tier~2 dispatches are labeled in the dashboard as
        \texttt{reconstructed (probabilistic)}, not \texttt{verified}.
  \item The dashboard tracks per-backend, per-goal-shape Tier~2
        reconstruction success rates over a rolling window.
  \item Tier~2 reconstruction failure messages distinguish ``backend
        did not produce a usable lemma list'' from ``lemma list
        produced but home system reconstruction failed''.
\end{itemize}

\paragraph{Demotion clause.} If implementation reveals that the
home-system reconstruction success rate from external lemma hints is
significantly worse than the Isabelle/\texttt{sledgehammer} baseline,
Tier~2 is to be demoted to a v1.5 milestone and v1 ships with Tiers
0, 1, 3 only.

\subsection{Tier 3: Proof trace}

\begin{lstlisting}[language=json]
"payload": {
  "trace_format": "alethe-2024" | "tstp-fof" | "tstp-thf" |
                  "lean-tactic-script" | "lfsc" | ...,
  "trace_data": <format-specific>,
  "trace_dialect_features": [ <feature>, ... ],
  "trace_annotations": <optional-natural-language>
}
\end{lstlisting}

\paragraph{Verifier behavior: strategy stack.} Tier~3 verification
proceeds through a strategy stack, in order, until one succeeds:

\begin{enumerate}[leftmargin=2em]
  \item \textbf{Native symbolic checker.} If the home system has a
        checker for this trace format (e.g., an Alethe-to-Lean trace
        replayer for Lean), use it. The checker consumes the trace and
        produces a home-system proof term step by step.
  \item \textbf{Cross-system replay.} Translate the trace into
        home-system tactics via a hand-written translation table for
        the format's rules. The resulting tactic script is run by the
        home-system and the kernel checks its output.
  \item \textbf{LLM-assisted reconstruction.} The trace, the goal, and
        any \texttt{trace\_annotations} are presented to a configured
        LLM with instructions to produce a home-system proof script.
        The script is kernel-checked. Failure is silent and proceeds to
        the next strategy.
  \item \textbf{Lemma extraction fallback.} If the trace can be summarized
        as a lemma list (e.g., the names of axioms used), produce a
        synthetic Tier~2 certificate from those names and re-verify at
        Tier~2.
\end{enumerate}

The user (or project configuration) controls which strategies are
enabled. LLM strategies are off by default in environments with no LLM
endpoint configured.

\paragraph{Trust expansion.} None. Each strategy produces a
home-system proof term that the kernel re-checks. LLM-proposed scripts
are sound iff the kernel accepts them.

\paragraph{Trace formats supported in v1.}

\begin{itemize}[leftmargin=2em]
  \item \texttt{alethe-2024}: SMT proof format produced by CVC5.
        Symbolic replay implemented for Lean.
  \item \texttt{tstp-fof}, \texttt{tstp-thf}: TPTP proof formats produced
        by Vampire. Symbolic replay implemented partially; LLM-assisted
        reconstruction available as fallback.
  \item \texttt{lean-tactic-script}: when the backend is an LLM-prover
        targeting Lean directly, the trace is the proposed script.
        Verified by running the script and checking the kernel result.
\end{itemize}

\subsection{Tier 4: Foundational certificate (reserved)}

Tier~4 certificates carry proofs in a foundational logical framework
(e.g., $\lambda\Pi$-calculus modulo) that the home system's kernel
accepts after a foundational embedding. Not implemented in v1.

\section{The dispatch protocol}
\label{sec:dispatch}

\subsection{Modes of dispatch}

The system distinguishes three modes of operation:

\paragraph{Fast path.} Native home-system automation (Lean's
\texttt{simp}, \texttt{aesop}, \texttt{decide}; Rocq's \texttt{auto},
\texttt{lia}, etc.) is unaffected by the brokerage. The brokerage is
invoked only via explicit \texttt{dispatch}-family tactics. Fast path
operates at sub-second latency; it is part of the home system, not
this specification.

\paragraph{Slow path (interactive).} The user invokes a
\texttt{dispatch} tactic. The plugin serializes the proof state, runs
the IR rewriter, hands off to the dispatcher, and awaits a result.
Results return asynchronously when possible: the IDE shows a pending
indicator on the tactic call until the dispatch resolves. State
versioning ensures that stale results do not get applied to mutated
proofs. Typical timeout: 10--60 seconds.

\paragraph{Build path.} On project compilation, dispatch calls re-run
with longer budgets and persistent caching. The build path is the
source of truth for CI; a project's CI passes iff every dispatch call
produces a kernel-checked proof term within budget. Cached certificates
from prior interactive runs are consulted before re-dispatch.

\subsection{The dispatch tactic API (Lean 4)}

\begin{lstlisting}[style=lean]
-- Generic dispatch with default strategy
example : 2 + 2 = 4 := by dispatch

-- Targeted dispatch
example : ... := by dispatch_smt (timeout := 10) (backend := cvc5)
example : ... := by dispatch_atp (backend := vampire)
example : ... := by dispatch_llm (model := default)

-- Dispatch with hints
example : ... := by
  dispatch (using := [Nat.succ_pos, foo_lemma])
           (tier := structural)
           (rewrite := default)
\end{lstlisting}

The Rocq API mirrors this via a \texttt{Dispatch} tactic with named
arguments; the surface syntax differs but the semantics are identical.

\subsection{Dispatch decision logic}

For v1, the dispatcher uses a rule-based decision procedure:

\begin{enumerate}[leftmargin=2em]
  \item If the user named a backend, use that one (after capability
        check).
  \item Otherwise, classify the goal by its
        \texttt{logic\_classification}:
        \begin{itemize}
          \item Pure linear arithmetic: SMT, prefer Tier~1 if a
                witness is feasible.
          \item First-order with equality, no arithmetic: ATP.
          \item Higher-order: ATP with HOL extensions, or LLM, or
                ITP if configured.
          \item Mixed: speculative parallel dispatch to top-2
                candidates.
        \end{itemize}
  \item Apply the project's preferences and exclusions (e.g., ``no LLM
        on this project'', ``prefer Z3 over CVC5'').
\end{enumerate}

A learned dispatcher is a v2 concern; v1 prioritizes predictability.

\subsection{Capability matching}
\label{sec:capability-matching}

Each backend adapter declares its capabilities as a structured manifest:

\begin{lstlisting}[language=json]
{
  "adapter": "cvc5",
  "logic_fragments": ["LIA", "LRA", "BV", "UF", "UFLIA", "UFLRA"],
  "type_constructions": ["primitive", "type_variable_via_specialization"],
  "max_order": "first_order",
  "tiers_produced": [1, 3],
  "trace_formats_produced": ["alethe-2024"],
  "preferred_rewrite_pipeline": [...]
}
\end{lstlisting}

Before invoking an adapter, the dispatcher checks the goal's
\texttt{logic\_classification} and the metadata-recorded type
constructions against the adapter's manifest. Mismatches cause the
dispatcher to skip the backend with reason
\texttt{logic\_out\_of\_fragment} or \texttt{type\_construction\_not\_supported}.

\subsection{Concurrent dispatch}

When the dispatcher chooses parallel speculative dispatch, multiple
adapters are invoked concurrently with the same (rewritten) IR. Result
handling policy:

\begin{itemize}[leftmargin=2em]
  \item \textbf{First valid certificate wins by default.} The first
        backend to return a certificate that the verifier accepts wins;
        other backends are cancelled.
  \item \textbf{Tier preference within a grace window.} If a higher-tier
        certificate (preferring Tier~1 over Tier~3 over Tier~2 over
        Tier~0) returns within a configured grace window after a
        first success, the higher-tier certificate is preferred.
        Default grace window: 2 seconds.
  \item \textbf{User override.} The user can specify
        \texttt{tier\_preference} in directives to force a different
        ordering.
\end{itemize}

\subsection{Failure handling}

\paragraph{Backend timeout.} Returns an \texttt{unknown} certificate
(envelope only, with \texttt{reason: "timeout"}). The dispatcher does
not silently fall through to another backend unless speculative dispatch
was configured.

\paragraph{Backend error.} Returns an \texttt{error} certificate with
backend-specific diagnostics. The user sees a clear message; the
dashboard records the error.

\paragraph{Verifier rejection.} The certificate is logged with
rejection reason; the dispatch is reported as failed. Speculative
parallel dispatches continue to run; a rejected certificate from one
backend does not affect others.

\paragraph{Reconstruction failure.} Reported per-strategy in the
strategy stack. The user sees a breakdown (e.g., ``Alethe-to-Lean
replay failed at step 47; LLM reconstruction also failed; falling
back to lemma extraction succeeded'').

\subsection{The dashboard}

The dashboard is a session-level UI element that surfaces dispatch
behavior over time. It tracks:

\begin{itemize}[leftmargin=2em]
  \item Per-call breakdown: which backends were invoked, which tiers
        were produced, which strategies succeeded, what timings were
        observed.
  \item Aggregate statistics: per-backend success rates over the
        session and the project, per-goal-shape Tier~2 reconstruction
        rates, LLM strategy success rates.
  \item Failure analysis: common failure modes, suggested remedies
        (e.g., ``Goals involving \texttt{LinearOrderedField} are
        succeeding 80\% on Z3 but 20\% on CVC5; consider
        \texttt{(backend := z3)}'').
\end{itemize}

The dashboard is the primary mechanism by which users develop
calibrated intuitions about dispatch behavior. Without it, the
brokerage operates as a black box; with it, the user has visibility
into the prover marketplace.

\section{The certificate verifier}
\label{sec:verifier}

\subsection{Architecture}

The verifier is a tier-dispatched function with per-tier modules. Each
module implements the following interface:

\begin{lstlisting}[style=plain]
trait TierVerifier {
  fn verify(
    cert: &Certificate,
    ir: &IR,
    rewrite_trace: &RewriteTrace
  ) -> VerificationResult;
}

enum VerificationResult {
  Verified(HomeSystemProofTerm),
  Rejected(RejectionReason),
  Indeterminate(IndeterminateReason),
}
\end{lstlisting}

A \texttt{Verified} result includes a home-system proof term of type
matching the (rewritten) goal. The \texttt{Rejected} variant carries a
structured rejection reason consumable by the dashboard.
\texttt{Indeterminate} is reserved for Tier~0 with pending trust
configuration.

\subsection{Pre-tier envelope verification}

Before tier-specific verification:

\begin{enumerate}[leftmargin=2em]
  \item Validate the envelope schema and version.
  \item Confirm \texttt{dispatch\_context\_hash} matches the actual
        IR document hash.
  \item Confirm \texttt{rewrite\_trace\_hash} matches the actual
        trace hash.
  \item Confirm the backend identity is recognized and not blocklisted.
  \item Validate library provenance: every referenced lemma's content
        hash matches a lemma available in the home system.
\end{enumerate}

Any failure here results in \texttt{Rejected} with reason
\texttt{envelope\_invalid}.

\subsection{Tier-specific verification}

\paragraph{Tier 0.} Consult trust policy. If accepted, produce a
proof term constructed via the home system's
\texttt{external\_oracle\_certificate} mechanism (which records the
trust expansion and the backend attestation in the proof term itself,
making the trust visible in any downstream consumer of the term).

\paragraph{Tier 1.} Dispatch to a per-witness-kind checker:

\begin{itemize}[leftmargin=2em]
  \item \texttt{farkas}: verify nonnegativity and arithmetic
        contradiction; produce a Lean/Rocq term using the home system's
        Farkas-witness lemma.
  \item \texttt{sat\_assignment}: evaluate the formula under the
        assignment; produce a term using the appropriate evaluation
        lemma.
  \item \texttt{sat\_unsat\_core}: run home-system propositional
        decision procedure on the core; produce a term.
\end{itemize}

Witness checkers are small, audited home-system tactics. Their
implementations are part of the home-system plugin.

\paragraph{Tier 2.} Invoke the home-system reconstruction tactic
indicated by \texttt{strategy\_hint}, with \texttt{lemmas\_used} as
hints, within the configured budget. Success iff the tactic produces
a kernel-checkable term.

\paragraph{Tier 3.} Invoke the strategy stack
(Section~\ref{sec:certificates}). Each strategy attempt is bounded by
a per-strategy timeout. Successful strategies short-circuit the stack;
failures fall through.

\subsection{Lifting (the inverse of refinement)}
\label{sec:lifting}

After tier-specific verification produces a proof term of the
\emph{rewritten} goal, the verifier hands off to the lifting phase.
Lifting inverts:

\begin{enumerate}[leftmargin=2em]
  \item Adapter refinement (using the certificate's
        \texttt{refinement\_record}).
  \item IR rewriting (using the rewrite trace).
\end{enumerate}

The output is a proof term of the \emph{original} home-system goal.

\paragraph{Adapter refinement inversion.} For each entry in the
\texttt{refinement\_record}'s \texttt{specializations} list, lifting
applies the inverse:

\begin{description}
  \item[\texttt{type\_specialization}.] Apply the named soundness
        witness (an embedding lemma) to lift the proof from the
        specialized type to the original type variable.
  \item[\texttt{method\_specialization}.] Replace the specialized
        operator with the typeclass method, using instance resolution.
  \item[\texttt{axiomatization}.] Discharge the introduced axioms
        using the home system's instance proofs.
  \item[\texttt{definition\_unfolding}.] Re-fold the definition
        (typically a no-op since unfolding preserves provability and
        the home system's elaborator handles the equivalence).
  \item[\texttt{logic\_encoding}.] Decode the encoding using a
        home-system decoder for that encoding family.
  \item[\texttt{rendering}.] No inversion needed; rendering is
        used only by adapters whose output is already in the home
        system's native syntax (LLM and ITP adapters).
\end{description}

\paragraph{IR rewrite inversion.} For each entry in the rewrite trace
(in reverse order), the corresponding pass's \texttt{inversion\_data}
is consulted to construct a wrapping that converts the proof of the
post-pass goal into a proof of the pre-pass goal.

\paragraph{Failure mode.} Lifting can fail when refinement record or
rewrite trace data is incomplete or inconsistent (a soundness-relevant
serializer or rewriter bug). On lifting failure, the verification
result is \texttt{Rejected} with reason \texttt{lifting\_failed} and
the partial proof term is discarded. The dispatcher does not retry
with the same backend; it reports the failure to the dashboard with
sufficient detail for offline diagnosis.

\subsection{Caching}

Successful certificates (post-lifting) are cached. The cache key is:

\[
(\textrm{goal\_hash}, \textrm{dispatch\_context\_hash}, \textrm{rewrite\_trace\_hash}, \textrm{backend\_id}, \textrm{backend\_version})
\]

Cache invalidation triggers:

\begin{itemize}[leftmargin=2em]
  \item Backend version change.
  \item Any referenced lemma's content hash change (detected via
        provenance hashes).
  \item Home system version change (defensive, conservative).
  \item Manual invalidation via dispatcher command.
\end{itemize}

Cache storage is project-local by default
(\texttt{.proof-broker-cache/} in the project root), with optional
shared-cache configurations for team or CI use.

\section{The reconstruction layer}
\label{sec:reconstruction}

\subsection{Per-home-system responsibilities}

Each home system's plugin provides:

\begin{enumerate}[leftmargin=2em]
  \item \textbf{Serializer.} \texttt{LocalProofState $\to$ IR}. Subject
        to the translation discipline (Section~\ref{sec:ir-discipline}).
  \item \textbf{Deserializer for verifier outputs.}
        \texttt{HomeSystemProofTerm $\to$ LocalTactic}. Inserts the
        verified proof term into the user's proof.
  \item \textbf{Translation tables.} Source-language types and symbols
        to IR shell entries, with provenance.
  \item \textbf{Witness checkers.} Per-witness-kind home-system tactics
        for Tier~1 verification.
  \item \textbf{Reconstruction tactics.} Lemma-hint consumers for
        Tier~2 verification (\texttt{aesop}-with-hints,
        \texttt{simp\_arith}-with-hints, etc.).
  \item \textbf{Trace consumers.} Per-format implementations for
        Tier~3 verification (Alethe replayer, TPTP replayer, etc.).
        Optional; missing consumers cause Tier~3 to fall back to
        LLM-assisted reconstruction or lemma extraction.
\end{enumerate}

\subsection{The translation contract}

The home-system plugin is bound by:

\begin{requirement}[Soundness contract]
For every IR document $D$ the serializer produces, and every
home-system proof term $\pi$ claiming $D$'s \texttt{goal} (in
shell form) under the rewriting and refinement records associated
with $D$, the lifting procedure must produce a home-system proof term
$\pi'$ such that the kernel accepts $\pi'$ at the type of the original
proof state's goal.
\end{requirement}

This is the central soundness obligation. Every other component
operates under the assumption that the home-system plugin satisfies
it. Failures of this contract are soundness bugs, not capability
limitations.

\subsection{Insertion into the proof}

Once lifting produces a proof term of the original goal type, the
plugin inserts it into the user's proof. For interactive use, this
manifests as the dispatch tactic succeeding and the goal being
discharged. For build-time use, the term is recorded in the cache and
re-used on next compilation.

\subsection{Reporting reconstruction in the proof term}

Where the home system's proof term language permits, the inserted term
is annotated with metadata indicating that it was produced by
dispatch. Examples:

\begin{itemize}[leftmargin=2em]
  \item For Lean 4, an inserted term may be wrapped in a
        \texttt{Proof.byDispatch} annotation that records the backend
        identity, tier, and certificate hash. This is invisible to the
        kernel but visible to project tooling.
  \item For terms accepted at Tier~0, the wrapping is mandatory and
        records the trust expansion explicitly.
\end{itemize}

This annotation is for auditability: a project team can run a tool to
list every dispatched proof in the codebase and the backends that
produced them.

\section{LLM integration points}
\label{sec:llm}

LLMs appear as first-class participants in three places, each gated by
configuration. In all cases, LLMs are proposers; the home system's
kernel is the disposer.

\subsection{LLM-as-backend}

A configured LLM endpoint is registered as a backend with its own
adapter. The adapter:

\begin{enumerate}[leftmargin=2em]
  \item Renders the IR as source-language syntax (Lean 4 surface
        syntax for the Lean home system, Rocq syntax for Rocq).
  \item Sends the rendering to the LLM endpoint with a prompt template
        instructing the model to produce a tactic proof.
  \item Receives the tactic proof and packages it as a Tier~3
        certificate with \texttt{trace\_format: "lean-tactic-script"}
        (or analogous for Rocq).
  \item Records its rendering choices in a refinement record with
        \texttt{kind: rendering}.
\end{enumerate}

The verifier processes the certificate by running the tactic script in
the home system. The kernel re-checks the resulting proof term.

\subsection{LLM-assisted Tier~3 reconstruction}

When a Tier~3 certificate from an SMT or ATP backend has a trace in a
format the home system cannot replay symbolically, the LLM-assisted
strategy may be tried:

\begin{enumerate}[leftmargin=2em]
  \item The trace, the goal, and any \texttt{trace\_annotations} are
        formatted into a prompt.
  \item The prompt asks the LLM to translate the trace into a
        home-system tactic proof.
  \item The resulting tactic proof is run and kernel-checked.
\end{enumerate}

This strategy is sound because the kernel checks the result; an LLM
that hallucinates simply produces a script that fails to type-check.

\subsection{LLM-assisted premise selection (deferred)}

For Tier ``rich'' context, an LLM can propose lemma slices that go
beyond symbol-based selection. Affects only what the backend sees;
does not affect soundness because the backend's certificate is still
kernel-checked. \emph{Deferred to v2.}

\subsection{Configuration}

LLM use is configured at the project level:

\begin{lstlisting}[language=json]
{
  "llm": {
    "enabled": true,
    "endpoint": "https://...",
    "model_identity": "<provider>/<model>/<version>",
    "max_tokens": 4096,
    "temperature": 0.0,
    "use_for": ["backend", "tier3_reconstruction"]
  }
}
\end{lstlisting}

LLM strategies are off by default in environments without LLM
configuration; their absence is not a failure mode.

\subsection{Cache invalidation for LLM strategies}

Cached certificates produced via LLM strategies are keyed in part on
\texttt{model\_identity}. Any reported model change invalidates
prior LLM-strategy cache entries. The home system additionally
re-verifies (kernel re-checks) cached LLM-reconstructed certificates
on load, even when model identity matches, as a defense against silent
model drift between cached identity reports.

\section{ITP-to-ITP dispatch}
\label{sec:itp-itp}

\subsection{The asymmetry}

Dispatch from one ITP to another (e.g., a Lean goal shipped to Isabelle
for proof) is qualitatively different from dispatch to SMT/ATP/LLM
backends. The asymmetry arises because the artifact returned by an ITP
is itself a foundational proof in another foundational logic, not a
backend-neutral proof artifact. Translating, say, an Isar proof to a
Lean proof term requires either:

\begin{enumerate}[leftmargin=2em]
  \item A foundational embedding from one logic to the other (the
        Tier~4 / Dedukti-style approach).
  \item Trusting the source ITP as an oracle (Tier~0).
  \item Reducing the proof to a lemma list and reconstructing in the
        target (Tier~2).
\end{enumerate}

For v1, options 2 and 3 are supported; option 1 is out of scope.

\subsection{Tier~0 ITP-to-ITP dispatch}

A configured ITP backend may be invoked at Tier~0 with explicit per-project
trust opt-in. The dispatcher renders the goal in the target ITP's
syntax (using the LLM-style rendering machinery), invokes the target
in batch mode (e.g., \texttt{isabelle build}), and accepts a
``proof succeeded'' result as a Tier~0 certificate. The trust expansion
is explicit and recorded in the proof term.

\subsection{Tier~2 ITP-to-ITP dispatch}

Many ITPs can report which lemmas were used in a proof: Isabelle's
\texttt{sledgehammer} returns proofs as \texttt{by (metis lemma1
lemma2)} and similar. The ITP adapter extracts this lemma list and
packages it as a Tier~2 certificate. Reconstruction in the home system
proceeds as for any Tier~2 certificate.

\paragraph{Limitation.} Tier~2 ITP-to-ITP dispatch only works when the
lemmas used are either: (a) hypotheses available in the home system's
local context, or (b) lemmas with cross-system equivalents identified
by content hash. Library lemmas without cross-system equivalents
cannot be used; the certificate is rejected with reason
\texttt{lemma\_not\_available\_in\_home\_system}.

\subsection{Future work}

Full Tier~3 reconstruction across ITPs (translating Isar proofs to
Lean tactic scripts symbolically) is a research problem deferred to
v2 or beyond. Tier~4 foundational interchange is explicitly
out of scope for the brokerage and best pursued by complementary
projects (Dedukti, Logipedia).

\section{Worked examples}
\label{sec:examples}

This section walks through three end-to-end dispatch examples,
illustrating how the architecture handles increasingly demanding
cases. Each example exercises different parts of the IR's metadata
machinery.

\subsection{Example 1: Linear arithmetic over a typeclass}
\label{sec:example1}

\paragraph{Source goal (Lean 4).}

\begin{lstlisting}[style=lean]
variable {alpha : Type*} [LinearOrderedCommRing alpha]

example (n m : alpha) (h1 : n + m = 10) (h2 : 0 <= n) (h3 : 0 <= m) :
    n <= 10 := by
  dispatch_smt
\end{lstlisting}

\paragraph{Walkthrough.}

\begin{enumerate}[leftmargin=2em]
  \item The serializer produces an IR document. Logic classification:
        \texttt{first\_order}, fragment \texttt{none-without-specialization}
        but with a \texttt{specialization\_targets} entry on the
        instance metadata pointing to LIA via the
        \texttt{embeds\_into:Int\_for\_universal\_LIA} tag.
  \item The dispatcher routes to the CVC5 adapter based on user
        directive (\texttt{dispatch\_smt}).
  \item The IR rewriter runs CVC5's preferred pipeline:
        propositional simplification, no quotient elimination needed,
        no aggressive unfolding (the operations are already at the
        typeclass-method level the adapter expects).
  \item The CVC5 adapter inspects the rewritten IR. The instance
        metadata's theory tag justifies type specialization
        $\alpha \to \mathbb{Z}$. The adapter records this in a
        refinement record.
  \item The adapter emits SMT-LIB:
        \begin{lstlisting}[style=plain]
(set-logic LIA)
(declare-const n Int)
(declare-const m Int)
(assert (= (+ n m) 10))
(assert (<= 0 n))
(assert (<= 0 m))
(assert (not (<= n 10)))
(check-sat)
(get-proof)
        \end{lstlisting}
  \item CVC5 returns \texttt{unsat} with an Alethe proof trace.
  \item The adapter packages the result as a Tier~3 certificate with
        the Alethe trace and the refinement record (recording the
        $\alpha \to \mathbb{Z}$ specialization with soundness witness
        \texttt{linear\_ordered\_comm\_ring\_lia\_embedding}).
  \item The verifier replays the Alethe trace against Lean's
        \texttt{Int} arithmetic, producing a Lean proof term of
        $\forall n\ m : \mathbb{Z}, n + m = 10 \to 0 \leq n \to 0 \leq m \to n \leq 10$.
  \item The lifting phase applies the embedding witness lemma to
        transport this proof to the typeclass-quantified version, then
        instantiates at the user's specific $\alpha$. The result is a
        Lean proof term of the original goal.
  \item The plugin inserts the term, the goal closes.
\end{enumerate}

\paragraph{Spec components exercised.} Type metadata (typeclass instance
with abstract signature, axioms, theory tags); definitional metadata
(typeclass methods); refinement record (\texttt{type\_specialization} and
\texttt{method\_specialization}); Tier~3 with symbolic Alethe replay;
lifting via embedding witness.

\subsection{Example 2: Higher-order goal with function composition}
\label{sec:example2}

\paragraph{Source goal (Lean 4).}

\begin{lstlisting}[style=lean]
example (f g : Nat -> Nat) (P : (Nat -> Nat) -> Prop)
    (h1 : forall h : Nat -> Nat, P h -> P (h .. h))
    (h2 : P f)
    (h3 : f = g) :
    P (g .. g) := by
  dispatch
\end{lstlisting}

(Here \texttt{..} stands for function composition; in real Lean
syntax this would be the dot-circle operator.)

\paragraph{Walkthrough.}

\begin{enumerate}[leftmargin=2em]
  \item The serializer produces an IR document with logic classification:
        \texttt{higher\_order}, features \texttt{function\_quantification}
        and \texttt{predicate\_argument}, no first-order fragment.
  \item Definitional metadata for \texttt{Function.comp} carries kind
        \texttt{defined\_function}, definitional equation, and
        \texttt{concept\_tag: function\_composition}.
  \item The dispatcher's capability matching excludes CVC5 (max\_order
        first\_order) and Z3 by default. Candidates: Vampire-HOL,
        LLM-prover.
  \item Speculative parallel dispatch is configured. Both adapters run.
  \item The Vampire-HOL adapter applies its preferred rewriting:
        definition unfolding for \texttt{function\_composition} (using
        the definitional equation from the metadata). The adapter
        emits TPTP-THF and invokes Vampire.
  \item The LLM adapter renders the goal as Lean syntax and queries the
        configured LLM endpoint for a tactic proof.
  \item In a typical run, the LLM returns
        \texttt{subst h3; exact h1 f h2} (or similar) within ~2 seconds.
        The verifier runs this in Lean; the kernel accepts the
        resulting term. Tier~3 certificate, format
        \texttt{lean-tactic-script}.
  \item Vampire's run, meanwhile, may take longer; the dispatcher's
        ``first valid certificate wins by default'' policy accepts the
        LLM result and cancels Vampire.
  \item No lifting is needed: the LLM adapter's refinement record
        contains only a \texttt{rendering} entry, which has no
        inversion. The proof term is the original goal type directly.
\end{enumerate}

\paragraph{Spec components exercised.} Logic classification (higher-order);
definitional metadata for defined functions; capability matching
(filtering out first-order-only backends); concurrent dispatch; LLM-as-backend
with kernel-checked output; trivial lifting (rendering is non-invertible
because no inversion is needed).

\subsection{Example 3: Quotient-type goal reduced to underlying type}
\label{sec:example3}

\paragraph{Source goal (Lean 4).}

\begin{lstlisting}[style=lean]
def myRel (n : Nat) : Int -> Int -> Prop := fun a b => n divides (a - b)
def MyZMod (n : Nat) : Type := Quot (myRel n)
def addMyZMod (n : Nat) : MyZMod n -> MyZMod n -> MyZMod n :=
  Quot.lift2 (fun a b => Quot.mk _ (a + b)) (...)

example (n : Nat) (hn : 0 < n) (x y z : MyZMod n) :
    addMyZMod n (addMyZMod n x y) z =
    addMyZMod n x (addMyZMod n y z) := by
  dispatch_smt
\end{lstlisting}

\paragraph{Walkthrough.}

\begin{enumerate}[leftmargin=2em]
  \item The serializer produces an IR document with type metadata for
        \texttt{MyZMod n} (kind \texttt{type\_constructor\_application},
        construction kind \texttt{quotient}, with
        \texttt{underlying\_type} $= \mathbb{Z}$, equivalence relation
        \texttt{myRel n}, elimination and equality principles).
        Definitional metadata for \texttt{addMyZMod} (kind
        \texttt{lifted\_to\_quotient}, underlying function
        \texttt{Int.add}).
  \item The dispatcher routes to CVC5.
  \item CVC5's preferred rewriting pipeline includes
        \texttt{quotient\_elimination}. The pass:
        \begin{itemize}
          \item Recognizes that the goal is an equality in
                \texttt{MyZMod n}.
          \item Applies the elimination principle three times
                (introducing representatives $a$, $b$, $c$).
          \item Applies the equality principle, reducing the goal to
                \texttt{myRel n ((a + b) + c) (a + (b + c))}.
          \item Unfolds \texttt{myRel n} using its definitional equation
                from definitional metadata.
          \item Records inversion data: the introduction points, the
                equivalence proofs needed to wrap the resulting proof.
        \end{itemize}
        After the pass, the goal is
        \texttt{n divides (((a + b) + c) - (a + (b + c)))}, which
        simplifies to \texttt{n divides 0}.
  \item Propositional simplification reduces this further; the result
        is in CVC5's reach.
  \item CVC5 returns an Alethe proof of \texttt{n divides 0}.
  \item The verifier replays the Alethe trace, producing a Lean proof
        of the underlying-type goal.
  \item Lifting first inverts the adapter refinement (definition unfolding,
        nothing else), then walks the rewrite trace in reverse: the
        \texttt{quotient\_elimination} pass's inversion data is consulted
        to wrap the proof in the appropriate \texttt{Quot.ind} and
        \texttt{Quot.sound} structure.
  \item The result is a Lean proof of the original quotient-type goal.
        The plugin inserts the term.
\end{enumerate}

\paragraph{Spec components exercised.} Type metadata for type
constructions (quotients); definitional metadata for lifted functions;
the IR rewriting layer with quotient elimination; rewrite trace and
lifting through trace inversion; tier~3 with symbolic replay; the
combined effect of rewriter inversion plus refinement inversion.

\subsection{What these examples establish}

The three examples taken together exercise every major architectural
component:

\begin{itemize}[leftmargin=2em]
  \item Each of the four metadata categories (logic classification, type
        metadata, definitional metadata, library provenance) appears
        in at least two examples.
  \item Each tier supported in v1 (1, 2, 3) is exercised; Tier~0 is
        used in ITP-to-ITP scenarios outside these examples.
  \item Three different backends (CVC5, LLM-prover, and via capability
        matching, the absence of fit for Z3) are exercised.
  \item The IR rewriting layer is exercised both with no rewrites
        (Example~2) and with non-trivial rewrites and inversion
        (Example~3).
  \item Lifting is exercised through type specialization (Example~1),
        rendering (Example~2, trivially), and rewrite-trace inversion
        (Example~3).
\end{itemize}

These examples are not exhaustive but they are believed to be
representative of the cases the v1 implementation must handle. The
spec authors recommend that any prototype implementation begin by
making these three examples work end-to-end before generalizing.

\appendix

\section{Registry of recognized patterns}
\label{app:registry}

The architecture relies on a shared registry of recognized patterns
that the IR's metadata refers to. The registry is itself a
specification artifact, separate from but coordinated with the IR
schema. This appendix describes its initial content for v1.

\subsection{Logical features}
\label{app:logic-features}

The \texttt{features\_used} list in
\texttt{logic\_classification} draws from:

\begin{itemize}[leftmargin=2em]
  \item \texttt{universal\_quantification\_over\_first\_order}
  \item \texttt{existential\_quantification\_over\_first\_order}
  \item \texttt{function\_quantification}
  \item \texttt{predicate\_argument}
  \item \texttt{equality\_at\_first\_order\_type}
  \item \texttt{equality\_at\_function\_type}
  \item \texttt{equality\_in\_quotient\_type}
  \item \texttt{equality\_in\_inductive\_type}
  \item \texttt{induction\_recursion}
  \item \texttt{dependent\_pair}
  \item \texttt{dependent\_function}
  \item \texttt{universe\_polymorphism}
\end{itemize}

\subsection{First-order fragments}
\label{app:fragments}

\begin{itemize}[leftmargin=2em]
  \item \texttt{LIA}: linear integer arithmetic
  \item \texttt{LRA}: linear real arithmetic
  \item \texttt{NIA}: nonlinear integer arithmetic
  \item \texttt{NRA}: nonlinear real arithmetic
  \item \texttt{BV}: bitvectors
  \item \texttt{UF}: uninterpreted functions with equality
  \item \texttt{UFLIA}, \texttt{UFLRA}: combined theories
  \item \texttt{ARRAY}: theory of arrays
  \item \texttt{FOL}: pure first-order logic with equality
  \item \texttt{HOL}: higher-order logic with equality
\end{itemize}

\subsection{Theory tags for type metadata}
\label{app:theory-tags}

Theory tags assert that universal sentences in a particular fragment
valid over a model are valid over instances of the type. Common tags:

\begin{itemize}[leftmargin=2em]
  \item \texttt{embeds\_into:Int\_for\_universal\_LIA}
  \item \texttt{embeds\_into:Rat\_for\_universal\_LRA}
  \item \texttt{embeds\_into:Real\_for\_universal\_LRA}
  \item \texttt{embeds\_into:BV32\_for\_universal\_BV}
  \item \texttt{is:finite}
  \item \texttt{is:decidable\_equality}
  \item \texttt{is:linearly\_ordered}
  \item \texttt{is:well\_founded}
\end{itemize}

\subsection{Axiom shape descriptors}
\label{app:axiom-shapes}

Used in \texttt{abstract\_axioms} entries for typeclasses. Each
descriptor is a string-encoded pattern that adapters can
expand into concrete axioms in their target logic.

\begin{itemize}[leftmargin=2em]
  \item \texttt{associativity:OP}: $\forall x, y, z.\ (x \cdot y) \cdot z = x \cdot (y \cdot z)$
  \item \texttt{commutativity:OP}: $\forall x, y.\ x \cdot y = y \cdot x$
  \item \texttt{identity\_left:OP:E}: $\forall x.\ e \cdot x = x$
  \item \texttt{identity\_right:OP:E}: $\forall x.\ x \cdot e = x$
  \item \texttt{inverse:OP:INV:E}: $\forall x.\ x \cdot \mathrm{inv}(x) = e$
  \item \texttt{distributivity\_left:OP1:OP2}:
        $\forall x, y, z.\ x \cdot (y + z) = x \cdot y + x \cdot z$
  \item \texttt{reflexivity:R}: $\forall x.\ x \mathrel{R} x$
  \item \texttt{transitivity:R}: $\forall x, y, z.\ x \mathrel{R} y \to y \mathrel{R} z \to x \mathrel{R} z$
  \item \texttt{symmetry:R}: $\forall x, y.\ x \mathrel{R} y \to y \mathrel{R} x$
  \item \texttt{antisymmetry:R}: $\forall x, y.\ x \mathrel{R} y \to y \mathrel{R} x \to x = y$
  \item \texttt{totality:R}: $\forall x, y.\ x \mathrel{R} y \vee y \mathrel{R} x$
  \item \texttt{monotonicity\_in\_arg:OP:R:N}: $\forall \vec{x}, y, z.\ y \mathrel{R} z \to \mathrm{op}(\dots, y \text{ at } N, \dots) \mathrel{R} \mathrm{op}(\dots, z \text{ at } N, \dots)$
  \item \texttt{positivity\_preservation:OP}: positive multiplication and similar
\end{itemize}

The registry is extensible; adapters may report unrecognized shapes
and the IR can fall back to providing explicit axiom statements in
shell form. Adding a new shape descriptor requires registry
maintainers to ensure that all participating adapters either
recognize it or use the explicit-axiom fallback.

\subsection{Concept tags for definitional metadata}
\label{app:concepts}

\begin{itemize}[leftmargin=2em]
  \item \texttt{function\_composition}
  \item \texttt{function\_identity}
  \item \texttt{function\_application}
  \item \texttt{integer\_addition}, \texttt{integer\_multiplication},
        \texttt{integer\_subtraction}, \texttt{integer\_negation},
        \texttt{integer\_modulo}, \texttt{integer\_division}
  \item \texttt{boolean\_and}, \texttt{boolean\_or}, \texttt{boolean\_not}
  \item \texttt{set\_membership}, \texttt{set\_union}, \texttt{set\_intersection}
  \item \texttt{list\_cons}, \texttt{list\_nil}, \texttt{list\_append},
        \texttt{list\_length}
\end{itemize}

The concept tags allow adapters to recognize standard mathematical
operations across home systems. Lean's \texttt{List.cons} and Rocq's
\texttt{cons} both carry \texttt{concept\_tag: list\_cons}, so an
adapter that knows about lists can handle both uniformly.

\subsection{Construction kinds for type metadata}
\label{app:construction-kinds}

For \texttt{type\_constructor\_application}:

\begin{itemize}[leftmargin=2em]
  \item \texttt{quotient}: see Example~3 (Section~\ref{sec:example3})
  \item \texttt{inductive}: ADT-like type with named constructors
  \item \texttt{coinductive}: dual of inductive, with destructors
  \item \texttt{sigma}: dependent pair $\Sigma x : A.\ B(x)$
  \item \texttt{pi}: dependent function $\Pi x : A.\ B(x)$ (rare in
        v1; usually appears in shell-level \texttt{Forall})
  \item \texttt{refinement}: $\{x : A \mid P(x)\}$
  \item \texttt{higher\_inductive}: HIT with point and path constructors
        (deferred to v2 for production support; metadata schema reserved)
\end{itemize}

\subsection{Adapter capability vocabulary}
\label{app:adapter-caps}

For backend adapter capability manifests:

\begin{itemize}[leftmargin=2em]
  \item \texttt{logic\_fragments}: from Section~\ref{app:fragments}.
  \item \texttt{type\_constructions}: which \texttt{construction\_kind}
        values the adapter handles, plus pseudo-values such as
        \texttt{type\_variable\_via\_specialization} (handles type
        variables only when a specialization is recorded).
  \item \texttt{max\_order}: \texttt{first\_order} or \texttt{higher\_order}.
  \item \texttt{tiers\_produced}: which certificate tiers the adapter
        can produce.
  \item \texttt{trace\_formats\_produced}: which trace formats the
        adapter emits at Tier~3.
\end{itemize}

\subsection{Maintenance}

The registry is versioned alongside the IR schema. Adding entries is a
non-breaking change; removing or renaming entries is breaking and
requires schema-version increments. The registry source of truth is
maintained as a separate document with cross-references from this
specification.

\section{Open questions and deferred work}
\label{app:open}

The following items are explicitly deferred or unresolved at the time
of this v1.0 specification.

\subsection{Resolved during specification development}

\begin{itemize}[leftmargin=2em]
  \item \emph{Tier~2 in v1?} Resolved: yes, framed as a fallback tier
        with explicit probabilistic-soundness language and dashboard
        surfacing (Section~\ref{sec:tier2}).
  \item \emph{Lean + Rocq scope?} Resolved: Lean 4 primary, Rocq as
        architectural probe with Tier~0+1 only
        (Section~\ref{sec:scope}).
  \item \emph{Rich context overcommitment?} Resolved: rich context
        narrowed to symbol-based selection plus manual augmentation
        plus provenance, no transitive closure
        (Section~\ref{sec:ir-tiers}).
  \item \emph{Per-backend serialization vs.\ general IR?} Resolved:
        general IR plus per-backend adapter refinement, with
        configurable IR rewriting between (Sections~\ref{sec:overview},
        \ref{sec:rewriting}).
  \item \emph{Metadata category convergence?} Resolved: four bounded
        categories (logic classification, type metadata, definitional
        metadata, library provenance) with conversion metadata
        deferred to v2.
\end{itemize}

\subsection{Open for v1 implementation}

\begin{enumerate}[leftmargin=2em]
  \item \textbf{Translation hints for typeclass-heavy contexts beyond
        the worked example.} The architecture has been validated
        against three worked examples (Section~\ref{sec:examples}). Real
        usage will surface typeclass configurations not anticipated by
        these examples. The specification authors recommend that the
        v1 implementation track translation-hints failures in the
        dashboard with sufficient detail to drive registry extensions.
  \item \textbf{Binary serialization format.} JSON is the canonical
        format for this specification. Production deployments should
        migrate to a binary format (CBOR or Protocol Buffers) for
        performance. The choice is deferred to implementation; the
        binary format must round-trip with JSON to preserve the
        canonical specification semantics.
  \item \textbf{Cache invalidation under home-system updates.} The
        current scheme keys on referenced lemma content hashes plus
        backend versions. Home-system kernel updates may change
        definitional unfolding behavior in subtle ways that don't show
        up in lemma hashes. The conservative defense is to invalidate
        on any home-system version change; a more refined scheme can
        be introduced if conservative invalidation proves too expensive.
  \item \textbf{Concurrent dispatch policy refinement.} The default
        ``first valid certificate wins, with a 2-second grace window
        for higher tiers'' is a starting point. Real usage data will
        reveal whether a different policy (e.g., always wait for
        higher-tier certificates, or never wait) is preferable. The
        configuration mechanism allows experimentation.
  \item \textbf{LLM model identity and drift.} The defense against
        silent model drift is mandatory re-verification on cache load
        for LLM-strategy certificates. This is conservative and
        expensive. Whether a less expensive approach (verification
        only on cache miss with explicit identity comparison) is
        acceptable depends on operational experience with model
        endpoints.
  \item \textbf{Opaque payload visibility in dispatch UX.} When a goal
        contains payloads opaque to the chosen backend, the user
        should be informed that whole portions of the goal are being
        treated as uninterpreted. The current spec calls for diagnostic
        breakdowns in the dispatch UX but does not fully specify the
        UI; this is left to home-system plugin implementations.
  \item \textbf{Tier 3 trace richness for multi-strategy
        reconstruction.} The optional \texttt{trace\_annotations} field
        allows backends to provide natural-language commentary on
        proof traces, which can help LLM-assisted reconstruction.
        Whether to require backends to provide this when feasible
        (vs.\ leaving it optional) is unresolved; the v1 spec leaves
        it optional.
\end{enumerate}

\subsection{Deferred to v2}

\begin{enumerate}[leftmargin=2em]
  \item \textbf{Production-quality Rocq integration.} The architectural
        probe in v1 validates that the IR is not Lean-shaped; v2 must
        bring Rocq up to feature parity with Lean.
  \item \textbf{Tier~4 foundational certificates.} Requires committing
        to a foundational embedding target (likely $\lambda\Pi$-modulo
        / Dedukti) and building per-home-system embeddings. A research
        project complementary to but separable from the brokerage.
  \item \textbf{ML-driven premise selection.} For Tier ``rich''
        context. Requires premise-selection model training, evaluation
        infrastructure, and integration with the IR rewriter as a
        pluggable selection module.
  \item \textbf{Conversion metadata.} For goals depending on
        non-normalized term distinctions, an explicit conversion
        metadata category is required as the fifth IR metadata
        category.
  \item \textbf{Full Tier~3 cross-ITP reconstruction.} Translating
        proof artifacts from one ITP to another at Tier~3. Currently
        constrained to Tier~0 (oracle) and Tier~2 (lemma list) for
        ITP-to-ITP dispatch.
  \item \textbf{Higher inductive types.} Construction kind reserved;
        production support requires more design work on the
        path-constructor metadata.
  \item \textbf{Mutual recursion in definitional metadata.} The schema
        currently treats each defined symbol independently. Symbols
        defined in mutual blocks need cross-references, deferred.
\end{enumerate}

\section{Glossary}
\label{app:glossary}

\begin{description}
  \item[Adapter.] Per-backend component that consumes (rewritten) IR
        and produces backend-native input; produces a refinement
        record with returned certificates.
  \item[Backend.] An automated reasoning system: SMT solver, ATP,
        LLM-prover, or another ITP, capable of attempting to discharge
        a proof obligation.
  \item[Certificate.] The artifact returned by a backend adapter on
        successful dispatch, organized into tiers.
  \item[Dispatch.] The process of routing a proof obligation from a
        home system through the brokerage to a backend.
  \item[Home system.] An interactive proof assistant that uses the
        brokerage. Lean 4 (primary) or Rocq (probe) in v1.
  \item[IR.] Intermediate Representation: the backend-independent
        document describing a proof obligation.
  \item[Lifting.] The process of converting a proof of the rewritten,
        refined goal back to a proof of the original goal, using
        rewrite trace and refinement record.
  \item[Refinement record.] A structured description of an adapter's
        translation choices, attached to the certificate.
  \item[Rewrite trace.] A log of IR-to-IR transformations performed by
        the rewriting layer, used for inversion during lifting.
  \item[Shell.] The logic-agnostic typed term calculus used in IR
        documents.
  \item[Tier.] Classification of certificates by strength of evidence
        (0: oracle, 1: witness, 2: lemma list, 3: trace, 4:
        foundational).
\end{description}

\end{document}
